An electron moves down the page through a magnetic field that points to the right.

\begin{center}
\begin{tikzpicture}[>=Stealth]
 \foreach \y in {0,0.8,...,3.2} \draw[->, gray, thick] (0,\y) -- (4.4,\y);
 \node[gray, right] at (4.4,3.2) {$\vec B$};
 \fill[blue!60!black] (2.2,2.75) circle (0.17); \node[white] at (2.2,2.75) {\scriptsize $-$};
 \node[right] at (2.4,2.85) {\small e$^-$};
 \draw[->, very thick] (2.2,2.5) -- (2.2,1.1) node[right] {$\vec v$};
\end{tikzpicture}
\end{center}

Elif finds the direction of the magnetic force on the electron in four steps:
\begin{itemize}
\item[1.] The electron is negative: the left hand.
\item[2.] Thumb along the velocity, down the page.
\item[3.] Index finger along the field, to the right.
\item[4.] The middle finger points into the page. The electron is negative, so the force points the other way: out of the page.
\end{itemize}

\begin{abcliste}
\abc One step is wrong. Which one, and why?
\abc In which direction does the magnetic force on the electron point?
\abc In which direction would the force point on a proton moving the same way?
\end{abcliste}
